\documentclass[12pt]{article}

\usepackage{amsmath}
\usepackage{amssymb}
\usepackage{amsthm}
\usepackage{geometry}
\usepackage{setspace}

\geometry{letterpaper, margin=1in}
\setstretch{1.15}

\newtheorem{theorem}{Theorem}
\newtheorem{definition}{Definition}
\newtheorem{remark}{Remark}

\begin{document}

\title{\textbf{A Dimensional Scaling Condition for \\ Mixed Gaussian Traces}}
\author{}
\date{}
\maketitle

\begin{abstract}
\noindent We study an analytic expression, the \textit{Mixed Gaussian Trace} $\mathcal{Z}(p, q)$, defined as the tensor product of the continuous standard Gaussian measure on $\mathbb{R}^p$ and the discrete heat trace on the flat torus $\mathbb{T}^q$, evaluated at the modular self-dual point $\tau = i$. By framing this trace within the theory of periods and utilizing the Chowla-Selberg formula, we show that $\mathcal{Z}(p, q)$ evaluates to a monomial in fractional powers of $\pi$ and $\Gamma(1/4)$ over $\mathbb{Q}(\sqrt{2})$. Invoking Chudnovsky's theorem on the algebraic independence of $\pi$ and $\Gamma(1/4)$, we establish that the exponent of $\pi$ exactly vanishes if and only if the dimensions satisfy the scaling relation $2p = 3q$. At the fundamental integer solution $(p=3, q=2)$, the evaluation lies in $\mathbb{Q}(\sqrt{2}, \Gamma(1/4))$.
\end{abstract}

\vspace{0.5cm}

\section*{1. The Mixed Spectral Trace}
We define a mixed-dimensional analytic object via spectral geometry. Let $\Delta_{\mathbb{T}^q}$ be the positive Laplace-Beltrami operator on the standard flat torus $\mathbb{T}^q = \mathbb{R}^q / \mathbb{Z}^q$. The eigenvalues of the Laplacian are $\lambda_n = 4\pi^2 \|n\|^2$ for $n \in \mathbb{Z}^q$. The trace of the heat kernel at spectral time $t > 0$ is given by the multidimensional Jacobi theta function:
\begin{equation}
    K_{\mathbb{T}^q}(t) = \operatorname{Tr}\left(e^{-t\Delta_{\mathbb{T}^q}}\right) = \sum_{n \in \mathbb{Z}^q} e^{-4\pi^2 \|n\|^2 t}
\end{equation}

We evaluate this trace at the specific spectral time $t = \frac{1}{4\pi}$, corresponding to the self-dual point $\tau = i$ under the modular transformation $\tau \mapsto -1/\tau$. This yields $K_{\mathbb{T}^q}\left(\frac{1}{4\pi}\right) = \theta_3(e^{-\pi})^q$. Concurrently, we evaluate the standard unnormalized Gaussian integral over $\mathbb{R}^p$.

\begin{definition}[The Mixed Gaussian Trace]
Let $\mathcal{Z}(p, q)$ be the product of the continuous Gaussian volume of $\mathbb{R}^p$ and the discrete heat trace on $\mathbb{T}^q$ at $t = \frac{1}{4\pi}$:
\begin{equation}
    \mathcal{Z}(p, q) = \left( \int_{\mathbb{R}^p} e^{-\|x\|^2/2} \, d^px \right) \times \operatorname{Tr}\left(e^{-\frac{1}{4\pi}\Delta_{\mathbb{T}^q}}\right)
\end{equation}
\end{definition}

\section*{2. CM Evaluation and the Period Ring}
The continuous integral trivially evaluates to $(2\pi)^{p/2}$. The self-dual torus evaluated at $\tau = i$ corresponds to the elliptic curve $E: y^2 = x^3 - x$, which has complex multiplication (CM) by the Gaussian integers $\mathbb{Z}[i]$.

By the Chowla-Selberg formula, the periods of CM elliptic curves evaluate analytically in terms of the Euler Gamma function at rational arguments. Specifically, the trace at $\tau=i$ resolves to:
\begin{equation}
    \theta_3(e^{-\pi}) = \frac{\Gamma(1/4)}{2^{1/2}\pi^{3/4}}
\end{equation}

Taking the product across arbitrary dimensions $(p,q)$ yields the exact closed form:
\begin{equation}
    \mathcal{Z}(p, q) = (2\pi)^{\frac{p}{2}} \cdot \left(\frac{\Gamma(1/4)}{2^{1/2}\pi^{3/4}}\right)^q = 2^{\frac{p-q}{2}} \cdot \pi^{\frac{2p-3q}{4}} \cdot \Gamma\left(\frac{1}{4}\right)^q \label{eq:trace}
\end{equation}

This expression establishes that $\mathcal{Z}(p,q)$ evaluates to a monomial in fractional powers of the algebraically independent periods $\pi$ and $\Gamma(1/4)$ over the real quadratic field $\mathbb{Q}(\sqrt{2})$.

\section*{3. Exponent Vanishing via Chudnovsky's Theorem}
In 1984, G. V. Chudnovsky proved that $\pi$ and $\Gamma(1/4)$ are algebraically independent over $\mathbb{Q}$. Consequently, no non-trivial polynomial relation exists between them over the algebraic numbers. This dictates that the fractional power of $\pi$ in Equation \ref{eq:trace} cannot be algebraically simplified or absorbed by any rational function of $\Gamma(1/4)$.

\begin{theorem}[Vanishing of the $\pi$-Exponent]
The exponent of $\pi$ in the Mixed Gaussian Trace $\mathcal{Z}(p,q)$ vanishes, so that $\mathcal{Z}(p,q) \in \mathbb{Q}(\sqrt{2}, \Gamma(1/4))$, if and only if the dimensional parameters satisfy the linear scaling relation:
\begin{equation}
    2p = 3q
\end{equation}
\end{theorem}

\section*{4. The (3,2) Dimensional Scaling Relation}
The minimal positive integer solution satisfying this scaling relation occurs at $p=3$ and $q=2$. When evaluated at this dimensional ratio, the exponent of $\pi$ exactly vanishes ($\pi^0 = 1$):
\begin{align}
    \mathcal{Z}(3, 2) &= 2^{\frac{3-2}{2}} \cdot \pi^0 \cdot \Gamma\left(\frac{1}{4}\right)^2 \nonumber \\
    \implies \quad &\boxed{ \mathcal{Z}(3, 2) = \sqrt{2} \cdot \Gamma\left(\frac{1}{4}\right)^2 \approx 18.5899\dots }
\end{align}

\begin{remark}
This strict dimensional scaling clarifies the structural origin of previously observed heuristic identities. Coupling a 1-dimensional continuous measure with a 2-dimensional discrete trace ($p=1, q=2$) yields an invariant strictly proportional to $\frac{\Gamma(1/4)^2}{\pi\sqrt{2}}$. The persistence of $\pi$ in the denominator is the algebraic consequence of a non-zero exponent of $\pi$, as the tuple $(1,2)$ fails the $2p=3q$ condition. Isolating $\mathcal{Z}(3,2)$ provides the exact arithmetic scaling where the exponent of $\pi$ vanishes, yielding an evaluation strictly within the subfield $\mathbb{Q}(\sqrt{2}, \Gamma(1/4))$.
\end{remark}

\end{document}
